%%%%%%%%%%%%%%%%%%%% author.tex %%%%%%%%%%%%%%%%%%%%%%%%%%%%%%%%%%%
%
% sample root file for your "contribution" to a contributed volume
%
% Use this file as a template for your own input.
%
%%%%%%%%%%%%%%%% Springer %%%%%%%%%%%%%%%%%%%%%%%%%%%%%%%%%%


% RECOMMENDED %%%%%%%%%%%%%%%%%%%%%%%%%%%%%%%%%%%%%%%%%%%%%%%%%%%
\documentclass[graybox]{svmult}

% choose options for [] as required from the list
% in the Reference Guide

\usepackage{type1cm}        % activate if the above 3 fonts are
                            % not available on your system
%
\usepackage{makeidx}         % allows index generation
\usepackage{graphicx}        % standard LaTeX graphics tool
                             % when including figure files
\usepackage{multicol}        % used for the two-column index
\usepackage[bottom]{footmisc}% places footnotes at page bottom


\usepackage{newtxtext}       %
\usepackage[varvw]{newtxmath}       % selects Times Roman as basic font

\usepackage[style=apa, natbib]{biblatex}
\addbibresource{references_chapter.bib}

\makeindex             % used for the subject index
                       % please use the style svind.ist with
                       % your makeindex program

%%%%%%%%%%%%%%%%%%%%%%%%%%%%%%%%%%%%%%%%%%%%%%%%%%%%%%%%%%%%%%%%%%%%%%%%%%%%%%%%%%%%%%%%%

\begin{document}

\title*{Formal Methods for Cognitive Security: Engagement, Gaze, and Cybersickness}
\titlerunning{Formal Methods}
\author{Author One\orcidID{0000-0000-0000-0000} and\\ Author Two\orcidID{0000-0000-0000-0000}}

\institute{Author One \at Institution, City, Country, \email{author.one@example.edu}
\and Author Two \at Institution, City, Country, \email{author.two@example.edu}}

\maketitle

\section{Introduction}
\label{sec:introduction}

Mixed-reality (MR) systems are moving rapidly out of the laboratory and into settings where the cost of a wrong decision is high. Pilots train in head-worn displays, surgeons rehearse procedures on virtual patients, search-and-rescue teams plan operations against synthetic terrain, and combat medics learn casualty care under simulated battlefield stress. In each of these settings the user is asked to attend to information that arrives through a rendering pipeline whose behavior is not perfectly stable. Frame rates fluctuate, network jitter delays head-tracking updates, scene complexity strains the renderer, and the user's own physiology drifts over a long session. When the pipeline degrades, the user's cognition does too: attention wanders, gaze loses its tight coupling to objects of interest, and discomfort accumulates until performance breaks down. Cognitive security in this setting is not about firewalls or cryptography. It is about whether the user remains in a cognitive state from which good decisions can still be made.

This chapter takes a particular position on how to make that question tractable. We argue that three measurable cognitive states---\emph{engagement}, \emph{gaze}, and \emph{cybersickness}---are concrete enough that formal and semi-formal methods can reason about them, and important enough that operational MR systems already need quantitative answers to questions about them. We use these three as the chapter's case studies. The emphasis throughout is on what kind of guarantee a given method actually provides: an empirical accuracy number on a held-out cohort is not the same thing as a confidence-bounded statistical certificate, which is not the same thing as a probabilistic posterior derived from a graphical model, which is not the same thing as a sound bound on a model's worst-case output. Treating ``formal methods'' as a single category obscures these distinctions; the chapter spends considerable effort separating them.

By the end of the chapter, readers should be able to: (i) define cognitive security in an MR context and explain why the user's engagement, gaze stability, and discomfort are decision-relevant signals; (ii) name the four rungs of an empirical-to-formal guarantee ladder and identify which rung is being exercised in a given paper or system; (iii) describe at least one concrete model and one concrete verification technique for each of the three case studies; (iv) explain why two cybersickness models with similar predictive accuracy can produce different operational claims because they are verified differently; and (v) articulate a small number of unresolved problems where the next generation of researchers can usefully push the field.

The chapter assumes that earlier chapters of the textbook have introduced cognitive security at a general level, including notions such as resilience versus mitigation, individual differences, and trust. We focus on the MR context and treat formal methods as a tool for resilient decision-making rather than as an abstract mathematical topic. Section~\ref{sec:foundations} introduces the verification toolbox we will use throughout the chapter. Section~\ref{sec:cases} works through the three case studies in turn and ends with a short cross-case synthesis. Section~\ref{sec:evaluation} returns to the toolbox to discuss what each kind of guarantee licenses an operational system to do. Section~\ref{sec:implications} connects the technical material back to the broader textbook themes of resilience, trust, and individual differences, and frames a small number of forward-looking research questions.

\section{Cognitive Security and the Formal Methods Toolbox}
\label{sec:foundations}

For the purposes of this chapter, cognitive security in MR is the property that user attention, comfort, and decision-making remain inside operationally safe regions despite measurable degradation of the rendering and sensing pipeline. ``Operationally safe'' is intentionally application-specific: a tactical-medicine trainee should not become so disoriented that they cannot complete a triage protocol; a search-and-rescue operator should orient to a target object within a deadline; a surgical resident should sustain visual attention long enough to recognize a critical anatomical landmark. In each case the safety property is a statement about the user's cognitive state, not about the system's bits or pixels.

A predictive model alone---no matter how accurate on a held-out test set---does not tell an operational system \emph{when to believe} its own prediction at decision time. Average accuracy describes typical behavior over the data we happened to collect; it does not tell us anything about an individual user, a particular minute of a particular session, or an input region we did not sample. Formal methods, in the sense used in this chapter, fill that gap. A formal method takes three inputs---a fixed trained model, a property to check (for example, ``predicted FMS stays below~4'', or ``predicted gaze enters the object-of-interest region within one second''), and a region of admissible inputs---and returns one of three outputs: a guarantee that the property holds throughout the region, a counterexample that demonstrates a failure, or a quantified probabilistic bound on how often the property holds. Different formal methods make different choices about which of those three outputs they aim for and what assumptions they require to produce it.

\subsection{The Empirical-to-Formal Guarantee Ladder}
\label{sec:ladder}

The verification techniques used in this chapter sit on a single ladder, ordered by how much the guarantee commits to. Each rung is exercised by at least one of the case studies in Section~\ref{sec:cases}, and the chapter returns to the ladder in Section~\ref{sec:evaluation} to discuss what each rung licenses an operational system to do.

\textbf{Rung~1, empirical evaluation.} A predictor is run on data it has not been trained on, and average error or accuracy is reported. Cross-validation, leave-one-user-out evaluation, and held-out cohorts all live here. The guarantee is statistical in a weak sense: ``on data drawn from a distribution similar to the test set, the model behaves on average like this.'' This rung is the baseline against which everything else is compared and is often the only rung reported in the literature. Its limitation is that it says nothing about a particular input or input region.

\textbf{Rung~2, statistical certification of a runtime predicate.} A binary predicate is defined over the model's output (for example, ``predicted gaze entered the target region within one second and stayed there for at least 0.4 seconds''), the predicate is evaluated empirically over many windows that fall inside a chosen operating region~$\Omega$, and an exact binomial confidence bound such as the Clopper--Pearson lower bound \citep{clopper1934} is computed. The output is a statement of the form: ``with confidence~$1-\alpha$, the predicate holds with probability at least~$p_{\text{LB}}$ whenever the system state lies in~$\Omega$.'' The gaze case study in Section~\ref{sec:gaze} uses this rung, with $\alpha=0.10$ and $p_{\text{LB}}\geq0.82$ inside an operating region defined by optical-flow and temporal-smoothness constraints \citep{gaze2025reliability}. The strength of the rung is that the guarantee is checkable at runtime: if measurable scene statistics drift outside~$\Omega$, the system knows the certificate no longer applies.

\textbf{Rung~3, probabilistic verification on a graphical model.} A Bayesian network is specified or learned over the variables of interest, and exact inference is performed to compute a posterior probability of an outcome under specified evidence. The chapter uses Dice \citep{holtzen2020dice}, a probabilistic-programming language in which the network can be encoded directly and queries of the form $P(X_q \mid C)$ can be answered exactly under both ``hard'' evidence (deterministic constraints) and ``soft'' evidence (uncertain observations). The cybersickness snapshot model in Section~\ref{sec:cybersickness} sits on this rung \citep{wu2025cybersickness_bn}. The strength of the rung is that the answer is a calibrated probability rather than a point estimate; the cost is that the guarantee is only as good as the assumed causal structure of the network.

\textbf{Rung~4, deterministic verification of a fixed model artifact.} A trained model is treated as a fixed mathematical object, and a sound algorithm proves that for every input in a bounded region the model's output stays within specified bounds---or it returns a counterexample. Two flavors appear in this chapter. For neural networks, abstract interpretation in the DeepPoly relational domain \citep{singh2019deeppoly}, scaled to recurrent architectures by Prover \citep{ryou2021prover}, propagates polyhedral approximations of the input region forward through the network. For tree ensembles, the VoRF framework \citep{tornblom2019vorf} encodes each decision tree as a set of linear constraints, formulates a mixed-integer linear program whose objective is the worst-case forest output, and solves it exactly. The engagement case study in Section~\ref{sec:engagement} uses the neural-network flavor; the extended-exposure cybersickness case in Section~\ref{sec:cybersickness} uses the tree-ensemble flavor. The strength of the rung is that the guarantee is sound: no input in the certified region can violate the property. The cost is conservatism: certified regions can be very narrow, especially for high-capacity neural networks.

\subsection{What Each Rung Commits To}
\label{sec:rungs}

The four rungs differ along three axes that are easy to confuse. The first is what is being verified. Rung~1 verifies an average; rung~2 verifies a probability of a binary predicate; rung~3 verifies a posterior distribution; rung~4 verifies a deterministic bound on the model output. The second is what assumptions the guarantee rests on. Rung~1 assumes the test data is exchangeable with the deployment data. Rung~2 additionally assumes the certified operating region~$\Omega$ is correctly identified at runtime. Rung~3 assumes the causal structure of the Bayesian network is correctly specified, which is itself learned heuristically. Rung~4 assumes only that the model artifact is exactly the one that was certified and that the input region is correctly bounded; the guarantee is otherwise model-conditional and sound. The third axis is computational price. Rung~1 is essentially free given the test data. Rung~2 is cheap because it is just empirical-frequency counting plus a confidence-bound formula. Rung~3 is exact but exponential in the worst case; the cybersickness BN works because its size is moderate and Dice's compilation is efficient on discrete graphs. Rung~4 is the most expensive and the most conservative: DeepPoly intervals tighten only on small regions, and VoRF MILP encodings can grow with the number of trees and the depth of the forest.

\begin{svgraybox}
\textbf{Notation reused throughout this chapter.} The Bayesian network used in the cybersickness case study factorizes as
\[
P(\mathbf{S}, \mathbf{H}, X_q) = \prod_i P(S_i)\,\prod_j P(H_j \mid \mathrm{Pa}(H_j))\,P(X_q \mid \mathrm{Pa}(X_q)),
\]
where $\mathbf{S} = (S_1,\dots,S_m)$ are MR system parameters (latency, brightness, spectral entropy, optical flow), $\mathbf{H} = (H_1,\dots,H_k)$ are physiological responses (heart rate, galvanic skin response, gaze and pupil signals), and $X_q$ is cybersickness severity \citep{wu2025cybersickness_bn}. The parent operator $\mathrm{Pa}(\cdot)$ encodes a hierarchical constraint that system parameters affect physiology, and physiology affects cybersickness, but not the other way around within one observation window.
\end{svgraybox}

A reader who keeps the four rungs straight will already extract more value from the case studies than a reader who treats ``formal methods'' as a single technique. The case studies themselves were chosen partly because they exercise different rungs on superficially similar problems.

\section{Case Studies: Engagement, Gaze, and Cybersickness}
\label{sec:cases}

Each of the three case studies below takes the same shape. We define the cognitive state of interest, describe the experimental setup, sketch the predictor briefly, report the headline empirical numbers, and then explain which rung of the ladder in Section~\ref{sec:ladder} the case actually exercises and what the resulting guarantee says. A short cross-case synthesis closes the section.

\subsection{Cybersickness}
\label{sec:cybersickness}

\textbf{What cybersickness is and why it belongs in cognitive security.}
Cybersickness is a cluster of symptoms including nausea, disorientation, oculomotor strain, and visual fatigue that arise during virtual or mixed reality exposure and can persist well after the session ends \citep{rebenitsch2016cybersickness}. The symptoms closely resemble motion sickness but are triggered not by physical movement in the real world but by a mismatch between what the visual system perceives inside the headset and what the vestibular and proprioceptive systems report about the body's actual position. When the rendering pipeline introduces latency, drops frames, or injects camera jitter, the visual scene shifts in ways the body's balance sensors did not anticipate. The brain interprets this conflict as a threat, and the resulting physiological response is what we call cybersickness.

Three broad categories of factors drive cybersickness in MR systems (adapted from \cite{}). First, adversarial manipulations are deliberate modifications to system parameters such as display brightness, latency, or field of view that exploit known perceptual vulnerabilities to induce disorientation. Second, system limitations including tracking drift, sensor noise, rendering instability, and network latency create inconsistencies between virtual stimuli and human perception, amplifying discomfort. Third, environmental scene factors such as high object density, rapid scene motion, abrupt contrast gradients, and depth perception conflicts overload sensory processing and worsen symptoms. In safety-critical MR deployments all three categories can be present simultaneously, making reliable cybersickness modeling both challenging and essential.

From a cognitive security perspective, cybersickness is not merely a comfort problem. It is a degradation of the cognitive state from which decisions must be made. A military trainee experiencing moderate disorientation is less likely to triage a casualty correctly. A surgical resident whose gaze is compromised by visual fatigue may miss a critical anatomical landmark. A search-and-rescue operator struggling with nausea may fail to recognize an anomaly in a synthetic terrain model. The loss of cognitive integrity is the real threat, and cybersickness is one of its causes. This places cybersickness squarely within the cognitive security framing introduced in Section~\ref{sec:foundations}: the property that user attention, comfort, and decision-making remain inside operationally safe regions despite measurable degradation of the sensing and rendering pipeline.

\textbf{How cybersickness is measured.}
Both models in this case study operationalize cybersickness severity using the Fast Motion Sickness scale (FMS, 0 to 10) \citep{keshavarz2011fms}, a single-item verbal rating that participants report by briefly interrupting the task to give a number. The scale is particularly well suited to repeated in-task assessment because it takes only a few seconds to collect and does not require a multi-item questionnaire. A score of zero means no discomfort at all, and a score of ten means the most severe discomfort imaginable; intermediate values track the progression from mild awareness through moderate nausea to severe incapacitation. A threshold of FMS greater than four is commonly used in the MR safety literature to define a moderate-symptom event, and a threshold of FMS greater than six defines a severe event.

The FMS scale captures the user's subjective experience, which is the operationally relevant quantity, but subjective reporting has known limits. Users differ in their baseline sensitivity to motion, their willingness to report discomfort early, and their interpretation of the numeric scale. These individual differences are not noise to be averaged away; they turn out to be among the strongest predictors of cumulative cybersickness risk, a point the extended-exposure model makes explicit.

\textbf{Two timescales, two models.}
Cybersickness presents itself across two operationally distinct timescales. At the snapshot timescale, risk depends on the current configuration of the rendering pipeline: a sudden increase in display brightness, a burst of camera jitter, or an abrupt rise in scene complexity can elevate symptoms within seconds. At the extended-exposure timescale, risk accumulates over a session as fatigue compounds initial discomfort, and individual susceptibility characteristics determined before the user ever puts on the headset account for much of the variance. These two timescales call for different modeling strategies.

The first model addressed here (adapted from \cite{}) addresses the snapshot problem using a Bayesian Network verified through exact probabilistic inference, sitting on Rung~3 of the guarantee ladder. The second model (adapted from \cite{}) addresses the extended-exposure problem using a Random Forest regressor certified through mixed-integer linear programming, sitting on Rung~4. The two are not competing approaches to the same question; they are complementary tools that answer different operational questions and produce different forms of safety evidence.

\textbf{Snapshot model: Bayesian Network with probabilistic verification (Rung~3).}
The first model builds a hierarchical Bayesian Network over system-level features, human physiological signals, and cybersickness severity (adapted from \cite{}). The organizing insight is a causal hierarchy: system parameters influence the body's physiological responses, and those physiological responses drive cybersickness severity. This hierarchy is encoded as a directed acyclic graph in which system parameter nodes have no incoming edges, physiological response nodes receive edges only from system parameters, and the FMS severity node receives edges only from physiological signals. The joint distribution over all variables factorizes as:

\begin{equation}
P(\mathbf{S}, \mathbf{H}, X_q) = \prod_i P(S_i)\,\prod_j P(H_j \mid \mathrm{Pa}(H_j))\,P(X_q \mid \mathrm{Pa}(X_q)),
\end{equation}

where $\mathbf{S} = (S_1, \dots, S_m)$ are system parameters including luminance, spectral entropy, histogram-of-oriented-gradient (HOG) scene complexity features, and temporal smoothness; $\mathbf{H} = (H_1, \dots, H_k)$ are physiological features including heart rate variability, galvanic skin response (GSR), gaze error angle, pupil diameter, and reaction time; and $X_q$ is cybersickness severity.

Structure learning used the Hill-Climbing algorithm with the hierarchical constraints enforced throughout, so that learned dependencies never violate the causal ordering. After learning, node reduction removed variables with no directed path to the FMS outcome node, resulting in a compact 12-node network whose inference complexity is tractable in real time. Continuous features were discretized into five bins using K-means clustering before learning, enabling exact rather than sampling-based inference.

\begin{figure}[t]
\centering
\includegraphics[width=0.7\linewidth]{figures/structure_learned.png}
\caption{Structure of the learned hierarchical Bayesian Network. System-level features (orange nodes) influence physiological responses (blue nodes), which in turn drive cybersickness severity measured by FMS (green node). Directed edges encode the learned probabilistic dependencies under the imposed causal hierarchy. Adapted from \cite{}.}
\label{fig:bn-structure}
\end{figure}

The network was evaluated on data from 32 participants performing a 15-minute virtual reality walking task. It achieved 62.12\% exact-label accuracy and 83.55\% within-one-level accuracy across five participant-level cross-validation folds, meaning that most errors were off by a single severity level rather than grossly misclassified. The learned network structure is shown in Figure~\ref{fig:bn-structure}; system parameters form the top tier, physiological signals form the middle tier, and FMS severity occupies the bottom.

The critical output of this model is not the accuracy number. It is the verification step. Because all variables are discrete and the network has been compiled into the Dice probabilistic programming language \citep{holtzen2020dice}, the posterior probability of any Boolean property can be computed exactly. Two kinds of evidence can be applied. Hard conditioning fixes a variable to a known range, for example requiring that luminance exceeds discretized level three. Soft conditioning specifies that a variable satisfies a condition with some stated probability, for example that scene complexity reaches level four with 90\% confidence. Both forms of evidence update the joint distribution and propagate through the learned network to yield a calibrated posterior probability over FMS severity.

Figure~\ref{fig:system-impact} shows how individual system parameters influence the posterior probability of elevated FMS. Luminance and HOG-based scene complexity features emerge as the parameters with the greatest relative effect, while temporal smoothness and spectral entropy show negligible standalone impact. This finding directly informs which parameters a safety-aware MR system should monitor or constrain.

\begin{figure}[t]
\centering
\includegraphics[width=0.6\linewidth]{figures/system_params_impact_seaborn.png}
\caption{Impact of individual system parameters on the posterior probability of elevated FMS. Luminance and HOG-based scene complexity show the greatest relative influence on cybersickness risk. Temporal smoothness and spectral entropy contribute minimally as standalone indicators. Adapted from \cite{}.}
\label{fig:system-impact}
\end{figure}

The most operationally significant verification result concerns the combined effect of multiple high-risk parameters. When luminance, HOG scene complexity, and galvanic skin response are all constrained to discretized levels of at least four simultaneously, the posterior probability of severe cybersickness (FMS greater than four) rises from a baseline of approximately 45\% to approximately 81\% (adapted from \cite{}). Figure~\ref{fig:fms-shift} shows this shift in the FMS distribution. In practical terms, roughly four out of five users would experience significant discomfort under those combined conditions. This is not a prediction about a particular user in a particular session; it is a calibrated posterior derived from the learned joint distribution, and it is exactly the form an adaptive system needs to decide when to intervene.

\begin{figure}[t]
\centering
\includegraphics[width=0.6\linewidth]{figures/fms_comparison.png}
\caption{Shift in the FMS severity distribution under elevated conditions. When luminance, scene complexity (HOG features), and galvanic skin response are simultaneously constrained to high discretized levels, the distribution shifts markedly toward severe cybersickness outcomes (dark blue) compared to the baseline distribution (light blue). Adapted from \cite{}.}
\label{fig:fms-shift}
\end{figure}

The framework also supports the inverse inference direction: given that a safety threshold must be met with high probability, what operational constraints does that impose on physiological signals? Enforcing that FMS remains at or below four with 95\% probability shifts the posterior distribution of galvanic skin response strongly toward lower levels, as shown in Figure~\ref{fig:gsr-constraint}. The result is a formally derived operational guideline: maintaining galvanic skin response within discrete levels one through three is a necessary condition for safe operation under those parameter settings. This reversal of inference direction, from a desired safety property back to required physiological bounds, is a capability that purely predictive models cannot provide.

\begin{figure}[t]
\centering
\includegraphics[width=0.6\linewidth]{figures/gsr_comparison.png}
\caption{Posterior distribution of galvanic skin response under a 95\% safety constraint on FMS. Enforcing low cybersickness probability eliminates high-risk GSR levels from the feasible range, providing a formally derived operational bound for real-time physiological monitoring. Adapted from \cite{}.}
\label{fig:gsr-constraint}
\end{figure}

The guarantee sits on Rung~3 of the ladder. It is an exact posterior probability computed under specified evidence on a graphical model whose causal structure encodes a domain assumption about how system parameters, physiology, and cybersickness relate. The strength of the guarantee is conditional on that structural assumption being correct, which is why Rung~3 is distinct from the deterministic bound that Rung~4 provides.

\textbf{Extended-exposure model: Random Forest with MILP verification (Rung~4).}
The second model (adapted from \cite{}) addresses a qualitatively different question: what happens to cybersickness across a 30-minute MR session when perturbations are sustained, fatigue accumulates, and individual susceptibility shapes much of the variance? The experimental setting was a Tactical Combat Casualty Care training scenario on a HoloLens~2. Thirty-two participants completed approximately 30-minute sessions that included controlled camera-jitter perturbations at three intensity levels (Low, Medium, High) applied after a four-minute baseline. FMS was sampled at twelve scheduled timepoints, yielding 285 usable participant-timepoint samples.

A Random Forest regressor \citep{breiman2001randomforests} of 600 trees was trained over 24 features spanning five families. Susceptibility and pre-perturbation baseline variables captured individual differences in vulnerability: hours of sleep, a motion-sickness questionnaire score, a migraine-history score, and baseline simulator sickness subscales. Exposure context features encoded session elapsed time and perturbation intensity level. Gaze and head dynamics summarized report-aligned windows of 3D gaze origin, gaze direction, head position, and derived angular kinematics. Visual scene statistics captured optical flow, contrast, spectral entropy, and spatial frequency from the rendering stream. Skin conductance level provided a single physiological arousal signal.

The model was evaluated under five-fold band-balanced cross-validation, which distributed samples proportionally across low, mid, and high FMS bands in each fold to ensure balanced coverage of the discomfort range. It achieved a mean absolute error of $1.005 \pm 0.040$ on the 0 to 10 FMS scale and explained more than 80\% of the variance in self-reported discomfort ($R^2 = 0.812 \pm 0.027$). At the moderate-symptom threshold of FMS greater than four, the model correctly identified 82.2\% of threshold-crossing episodes with 84.1\% specificity. At the severe-symptom threshold of FMS greater than six, specificity rose to 98.0\%, indicating very few false alarms for extreme discomfort detection.

Temporal analysis revealed an important structure in how risk evolves across a session. The probability of predicted FMS exceeding four remained below 35\% through the first 16 minutes of the session and then rose sharply, crossing 50\% near minutes 28 to 30. This trajectory defines a natural intervention window: the system can operate with reasonable confidence that symptoms remain controlled through mid-session, but late-session risk requires proactive attention well before symptoms become severe.

The verification step applied the VoRF framework \citep{tornblom2019vorf} to the exported Random Forest. VoRF encodes every decision tree in the ensemble as a set of mixed-integer linear constraints and formulates an optimization problem whose objective is the worst-case forest output over a bounded input region. Solving this problem exactly yields a model-conditional deterministic certificate.

The certified region that satisfies FMS at or below four requires all of the following conditions to hold: hours of sleep of at least seven, a motion-sickness questionnaire score of at most two, a pre-session FMS reading at minute four of at most two, a migraine-history score of at most eleven, and a perturbation intensity of at most moderate. When all five conditions are met, no admissible combination of the remaining features can drive the model output above four. This is a sound guarantee about the model's behavior over every point in that region, not just on average. Outside this region, no formal certificate applies. Regions outside the certificate are not declared unsafe; the absence of a certificate simply means the formal guarantee no longer holds.

The feature importance structure carries a key message for cognitive security. Susceptibility and pre-session baseline variables together account for approximately 62\% of impurity-based importance across the ensemble, while runtime gaze and visual scene features contribute roughly 10\%. Session elapsed time contributes an additional meaningful share. Who the user is matters more than what the rendering pipeline is doing at any given moment, at least for predicting cumulative FMS across a 30-minute session. A safety system that relies only on real-time scene statistics will systematically underestimate vulnerability in high-susceptibility users, precisely the population that most needs early intervention.

\textbf{Two models, two operational shapes, one cognitive security framework.}
Keeping both models in a single case study is not redundant; it is pedagogically essential because the two models produce different forms of safety evidence suited to different points in the deployment pipeline.

The Bayesian Network model produces a calibrated posterior probability over FMS severity given observed or hypothesized evidence about current system and physiological conditions. This is a risk-quantification tool designed for in-session use. When the posterior probability of FMS exceeding four rises above an application-specific threshold, the system can trigger an adaptive response: reducing display brightness, simplifying scene content, recommending a brief rest, or in the extreme case initiating a session abort. The formal guarantee is probabilistic and conditioned on the causal structure of the network. It is updated continuously as new evidence arrives, which makes it well matched to the dynamic character of in-session risk.

The Random Forest model with VoRF certification produces a deterministic bound on the model output over a pre-specified region of admissible inputs. This is a screening tool designed for pre-session use. Before the session begins, the system can check whether the user's susceptibility profile and planned exposure intensity fall within the certified safe region. If they do, the model provides a sound guarantee that predicted FMS cannot reach the moderate-symptom threshold given the other features remain within their observed admissible ranges. This guarantee is binary: the user either satisfies the certified conditions or does not, and no probability threshold needs to be calibrated.

The two tools sit on different rungs of the guarantee ladder for reasons that reflect their operational purposes. A posterior probability (Rung~3) is the right output format for an in-session risk-thresholded intervention because it quantifies uncertainty and can be updated as evidence changes. A deterministic worst-case bound (Rung~4) is the right output format for a pre-session screening decision because it eliminates uncertainty about whether the threshold can be exceeded: the model is either certified or it is not, and the certificate is valid for every admissible input in the region. Neither substitutes for the other. A complete cognitive-security framework for MR cybersickness needs both: screening to select users into safe regions before exposure begins, and continuous probabilistic monitoring to detect when cumulative risk is approaching the intervention threshold during the session.

\subsection{Engagement}
\label{sec:engagement}

\textbf{What engagement is and why it belongs in cognitive security.}
Engagement, in the MR setting used here, is a latent cognitive state that reflects sustained attention to task-relevant virtual content \citep{engagement2026mr}. The construct is deliberately narrower than several related ones with which it is sometimes confused. \emph{Presence} captures the subjective feeling of ``being there'' in a virtual environment. \emph{Immersion} refers to properties of the system and display, such as field of view and rendering fidelity, that determine the perceptual envelope of the experience. \emph{Workload} characterizes cognitive demand, regardless of whether the user is actually attending to the relevant content. Engagement, by contrast, is the property that the user remains behaviorally and cognitively coupled to the task even when the rendering and interaction pipeline are being perturbed. It is influenced jointly by human-level features (gaze patterns, head dynamics, prior MR familiarity) and system-level parameters (latency, resolution, luminance), and by their temporal evolution within a session.

From a cognitive security standpoint, engagement is not a user-experience concern; it is the operational state from which decisions get made. A pilot trainee whose attention drifts during a synthetic instrument approach is not learning the procedure. A surgical resident who disengages from the virtual patient is not building the perceptual habits the simulation is designed to instill. A search-and-rescue operator who stops orienting to objects of interest is not extracting the situational awareness the MR interface is meant to deliver. This case study treats engagement as a measurable cognitive-security state and asks how MR systems can both estimate it under stress and provide formal guarantees about the conditions under which it is preserved.

A useful framing for those guarantees is the \emph{cognitive attack}: an adversarial or system-level perturbation of the MR interaction loop that degrades attention, situational awareness, or task execution \citep{engagement2026mr}. We focus on induced latency as a concrete instance because it is operationally relevant in MR, directly tied to rendering and interaction fidelity, and easily injected in controlled study. Other cognitive attacks (luminance shifts, jitter, tracking dropouts) follow the same template.

\textbf{How engagement is measured.}
Engagement is reported on a 1--7 Likert scale, sampled once per minute through an in-headset prompt. A score of~1 denotes complete disengagement from the MR content and~7 denotes full engagement. The single-item, minute-level format is well suited to repeated in-task assessment because it does not interrupt the task for long, but it inherits the usual limits of subjective self-report: users differ in their reporting style, the label distribution is uneven across both individuals and engagement levels, and there is no objective physiological ground truth against which the rating can be cross-validated. The sparsity and ordinality of the label shape two downstream choices that the predictor below makes explicit: contrastive pretraining as a way to learn structure from weak supervision, and a per-engagement-level evaluation rather than a single pooled error metric.

\textbf{The cognitive-attack setup.}
The study \citep{engagement2026mr} ran on a Microsoft HoloLens~2 in a controlled room-scale setup. Eleven male college students aged 25--32 each completed a single ten-minute mixed-reality hostage-rescue scenario containing hostages, enemies, fires, medical kits, and weapons. Participants were free to move, inspect the scene, and interact with virtual elements, but no explicit score or time-critical objective was imposed. The session was divided into ten one-minute intervals. In each interval the system was perturbed with a step change in rendering latency, drawn pseudo-randomly from three intensities: low (45--60~FPS), medium (20--45~FPS), and high (5--20~FPS). The next latency level was applied immediately after the participant submitted the engagement rating for the previous interval. Throughout each session, high-frequency multimodal telemetry was recorded continuously: binocular eye tracking (3D eye positions, gaze directions, and ocular rotations), gaze-target encodings identifying the object category being attended to, and 6-DoF head pose. The system stream contributed a latency-perturbation indicator and nine video-derived scene descriptors --- optical-flow magnitude, histogram-of-oriented-gradients features, edge intensity, temporal smoothness, brightness flicker, spectral entropy, spatial frequency, luminance, and contrast. These streams were sampled at their native rates and later aligned to the minute-wise engagement labels. The resulting dataset is small but tightly controlled, with both human and system signals densely sampled across all three perturbation conditions.

\textbf{Predictor: a two-phase contrastive encoder with a regression head.}
The predictor is a two-phase architecture designed for this regime of sparse ordinal labels and entangled multimodal inputs \citep{engagement2026mr}. In the first phase, an LSTM encoder \citep{hochreiter1997lstm} processes the variable-length multimodal trajectory $X = \{x_1, \ldots, x_T\}$, where each $x_t$ stacks the human and system features at time $t$, and maps it through masked average pooling and an L2-normalized projection into a 128-dimensional embedding $z$. The encoder is trained with a \emph{soft InfoNCE} objective \citep{khosla2020supervised} that uses the ordinal proximity of engagement labels as supervision: for two trajectories $i$ and $j$ with reported scores $y^{(i)}, y^{(j)} \in \{1,\dots,7\}$, the similarity weight is $\omega_{ij} \propto \exp(-|y^{(i)} - y^{(j)}|/\tau_s)$, so trajectories with similar engagement are pulled together in the embedding space and dissimilar ones pushed apart. With cosine similarity $\mathrm{Sim}(z_i, z_j)$ between embeddings, the contrastive loss takes the form
\begin{equation}
\mathcal{L}_{\text{con}} \;=\; -\sum_{i} \omega_{ij}\,\log\!\frac{\exp\!\bigl(\mathrm{Sim}(z_i, z_j)/\tau\bigr)}{\sum_{k \neq i} \exp\!\bigl(\mathrm{Sim}(z_i, z_k)/\tau\bigr)},
\end{equation}
where $\tau > 0$ is a separate temperature for the embedding space. In the second phase the encoder is frozen and a two-layer feedforward regression head $\hat{y} = W_2\,\mathrm{ReLU}(W_1 z + b_1) + b_2$ maps each embedding to a continuous engagement score in $[1, 7]$. The training objective combines mean-squared error with a variance regularizer over per-engagement-level errors, so the model does not concentrate accuracy on mid-scale ratings, and stratified mini-batch sampling keeps rare engagement levels represented in every batch. Figure~\ref{fig:engagement-overview} summarizes the full pipeline from multimodal MR telemetry through the contrastive encoder and regression head to the downstream verification step.

\begin{figure}[t]
\centering
\includegraphics[width=\linewidth]{figures/engagement_teaser.pdf}
\caption{Overview of the engagement case study. Multimodal MR signals (eye tracking, head motion, scene dynamics, and system latency) feed a contrastive LSTM encoder that learns engagement-aligned embeddings under a cognitive attack (here, latency perturbation). A calibrated regressor maps embeddings to a continuous engagement score, and the trained model is analyzed with formal verification to derive conservative operating envelopes for controllable MR parameters. Adapted from \citet{engagement2026mr}.}
\label{fig:engagement-overview}
\end{figure}

\textbf{Empirical results (Rung~1).}
Under label-stratified four-fold cross-validation, the model achieves a mean absolute error of $0.94 \pm 0.24$ on the 1--7 engagement scale, ahead of pooled-feature baselines on the same splits (Ridge 1.09, SVR 1.24, XGBoost 0.96). On three users excluded entirely from training and cross-validation, the model reaches MAE $0.84$ across 1{,}200 holdout samples --- \emph{better} than the cross-validation average. That gap is the cognitive-security signal that matters here: the contrastive embedding suppressed user-identity nuisance variation rather than memorizing the training participants, which is exactly the property a deployment-time predictor needs.

Table~\ref{tab:engagement-importance} reports the top ten features by permutation importance, the increase in MAE when each feature is independently shuffled at test time \citep{engagement2026mr}. The injected latency intensity is by far the strongest single feature, as expected for a cognitive-attack study, but it does not stand alone. Scene descriptors (spatial frequency, temporal smoothness, luminance) and gaze direction features (right gaze $z$, right gaze $y$) each contribute meaningfully, so the predictor is reading both the attack channel and how the user visually and behaviorally responds to it. This matters for cognitive security: a system that watches only its own perturbations is blind to whether those perturbations actually disrupted the user, whereas a system that watches the user's response alone misses the controllable system parameters that triggered the disruption in the first place.

\begin{table}[t]
\centering
\caption{Top ten features for the engagement regressor by permutation importance. The importance score is the increase in MAE when the feature is permuted at test time; larger values indicate features the model relies on more heavily. Adapted from \citet{engagement2026mr}.}
\label{tab:engagement-importance}
\footnotesize
\setlength{\tabcolsep}{8pt}
\begin{tabular}{lr}
\hline
\textbf{Feature} & \textbf{Importance (MAE Increase)}\\
\hline
AttackIntensity              & 1.0324\\
Feature\_spatial\_frequency    & 0.4454\\
Feature\_temporal\_smoothness  & 0.4145\\
RightGazeDirectionZ          & 0.1865\\
RightGazeDirectionY          & 0.1688\\
Feature\_luminance            & 0.1216\\
ObjectCount\_encoded          & 0.0328\\
Feature\_spectral\_entropy     & 0.0057\\
HeadRotationY                & 0.0055\\
LeftEyePositionZ             & 0.0045\\
\hline
\end{tabular}
\end{table}

\textbf{Verification (Rung~4).}
The empirical numbers above tell us how the model behaves \emph{on average}. They do not tell an MR system designer whether a given operating envelope of system parameters is \emph{guaranteed} to keep engagement above a desired level. That is a Rung~4 question, and it is the formal-methods contribution of this case study. Given the trained engagement regressor $f$ and a bounded region $I$ of MR system-parameter values, the verification question is whether the model's prediction stays above an engagement threshold $\tau$ for every input in that region:
\begin{equation}
\forall\, x \in I:\; f(x) \;\geq\; \tau.
\end{equation}
We report two natural thresholds. The first, $\tau = 4$, marks the midpoint of the 1--7 scale and separates distracted from attentive states. The second, $\tau = 5$, the second-highest score actually observed in the study, represents strong engagement. Verifying the predicate at both thresholds therefore captures both an acceptable and a higher standard of cognitive coupling. The verifier is Prover \citep{ryou2021prover}, which discharges the property by propagating DeepPoly relational abstractions \citep{singh2019deeppoly} through the unrolled LSTM and its nonlinear gates. If the property holds across the polyhedral over-approximation of the region, the certificate is sound; if a counterexample is found, the property fails. Table~\ref{tab:engagement-verification} reports the resulting certified intervals over the four most influential controllable system parameters.

\begin{table}[t]
\centering
\caption{Verified safety intervals for the engagement regressor using DeepPoly \citep{singh2019deeppoly} and Prover \citep{ryou2021prover}. Each cell reports a bounded range of an MR system parameter for which the predicted engagement is provably at least $\tau$ over every admissible input in the region. Intervals are sound but tight, reflecting the worst-case combination of stochastic user behavior and controllable system parameters. Adapted from \citet{engagement2026mr}.}
\label{tab:engagement-verification}
\footnotesize
\setlength{\tabcolsep}{8pt}
\begin{tabular}{lcc}
\hline
\textbf{System Parameter} & \textbf{Engagement $\geq 4$} & \textbf{Engagement $\geq 5$}\\
\hline
Luminance           & $[110.2,\,111.0]$ & $[109.8,\,110.6]$\\
Optical Flow        & $[0.021,\,0.024]$ & $[0.019,\,0.022]$\\
Spatial Frequency   & $[943.5,\,944.2]$ & $[944.0,\,945.2]$\\
Temporal Smoothness & $[7.1,\,8.0]$     & ---\\
Spectral Entropy    & ---               & $[7.01,\,7.05]$\\
\hline
\end{tabular}
\end{table}

The certified intervals are extremely tight: luminance is confined to a window of less than a single unit, optical flow to a few thousandths. This conservatism is not a bug in DeepPoly. It is a structural consequence of asking a deterministic worst-case question over a region in which much of the input is not controllable. The engagement model receives 40 inputs at every timestep, mixing controllable system parameters (latency, luminance, scene complexity) with uncontrolled user features (gaze stability, blink frequency, head motion). To certify the predicate for \emph{all} admissible inputs in a region, the verifier must absorb the worst-case user behavior as well as the worst-case system perturbation, and that worst case is what shrinks the certified slice of the input space.

The right reading of Table~\ref{tab:engagement-verification} is therefore not that deterministic NN verification is broken. It is that soundness has a price, and the price is paid in the size of the certified region. This is the chapter's central pedagogical move: the gaze case study in Section~\ref{sec:gaze} trades soundness for a statistical certificate over a runtime-checkable domain (Rung~2), and the cybersickness Bayesian network in Section~\ref{sec:cybersickness} uses probabilistic verification on a learned graphical model (Rung~3). The engagement case shows what soundness costs in practice on a high-capacity neural network; the next two cases show what techniques fill the gap when sound deterministic verification produces certified ranges that are too narrow to deploy.

\subsection{Gaze}
\label{sec:gaze}

\textbf{Why gaze.} Gaze is interesting in cognitive security for a reason that is easy to state empirically: in this dataset, gaze dwell duration on the target object correlates with mission-completion latency at $r=0.504$, $p=0.0053$ \citep{gaze2025reliability}. Users who maintain deliberate fixation on the right thing finish their tasks faster, which is exactly the pattern one would expect if gaze stability were a usable proxy for engaged, task-coupled attention. Gaze also has the operational advantage of being measurable in real time, so a system can act on a gaze certificate on the order of hundreds of milliseconds.

\textbf{Setup.} Data come from a tactical search-and-rescue MR testbed evaluated under leave-one-user-out cross-validation across thirty users. The display resolution is $2304 \times 2064$ and the nominal eye-tracking rate is approximately 60~Hz, with substantial irregularity introduced by rendering variability. Two attention properties are defined. The first is \emph{timely orientation}: predicted gaze should enter an object-of-interest region within $T_{\text{resp}}=1.0$~second of its appearance. The second is \emph{sustained attention}: gaze should remain in the region for at least $D=0.4$~seconds, the threshold conventionally separating deliberate fixation from a transient glance.

\textbf{Predictor.} The predictor is a Time-Aware LSTM \citep{baytas2017tlstm} that explicitly models the elapsed time $\Delta t$ between observations. Its hidden state is attenuated by a learned exponential decay $\gamma_t = \exp(-W_\gamma \Delta t_t)$ before the recurrent update, which prevents stale information from dominating after large gaps in the data stream. Sparse asynchronous events and object detections are folded in via a last-observation-carried-forward scheme with their own staleness decay, and the resulting fused features feed a regression head that predicts normalized gaze coordinates. On a 100~ms horizon the model achieves normalized RMSE of $0.059 \pm 0.015$ ($128.8 \pm 34.5$ pixels), a 27\% reduction relative to a standard LSTM and a much larger reduction relative to a persistence baseline. Cross-user generalization is tight: the interquartile range of normalized RMSE across the 30~LOUO folds is only 0.023.

\textbf{Verification (rung~2).} The verification question is not whether the network's outputs are bounded---it is whether the two attention properties hold often enough to be operationally trustworthy. The framework defines a domain~$\Omega$ over runtime-measurable scene statistics (optical flow and temporal smoothness), evaluates the binary predicate ``the property held in this window'' over many windows that fall inside~$\Omega$, and computes the Clopper--Pearson lower bound on the success probability. The published primary certificate states that whenever optical flow is in $[0.02, 0.07]$ and temporal smoothness is in $[20, 55]$, the orientation-and-fixation property holds with probability $p_{\text{LB}}\geq 0.82$ at 90\% confidence \citep{gaze2025reliability}. A separate bounded-divergence certificate gives stronger numbers ($p_{\text{LB}}\geq 0.998$ in the lowest-perturbation FPS bin, decaying to $0.970$ in the highest), and a window-length sensitivity analysis confirms that the certificate is not brittle to the choice of $T_{\text{resp}}$.

The framework explicitly does \emph{not} claim formal NN verification. The authors are precise about this: the guarantee is a statistical certificate of a runtime predicate over an empirically defined domain, not a deductive proof about the network's outputs \citep{gaze2025reliability}. Inside~$\Omega$ the system is licensed to act on gaze; outside~$\Omega$ the certificate does not apply and the system should fall back to mitigation strategies such as target enlargement or rendering stabilization. This is exactly the operational shape that rung~2 supports.

\begin{figure}[t]
\centering
%\includegraphics[width=0.9\textwidth]{fig_gaze_arch}
\rule{0.9\textwidth}{0.5pt}\\[0.5em]
\textit{[Placeholder: Time-Aware LSTM with learned exponential decay and sparse-context fusion.]}
\caption{Time-Aware LSTM architecture for gaze prediction. Hidden state is attenuated by a learned exponential decay before each recurrent update; sparse events and object detections are fused via decayed last-observation carrying. Adapted from \citet{gaze2025reliability}.}
\label{fig:gaze-arch}
\end{figure}

\begin{figure}[t]
\centering
%\includegraphics[width=0.7\textwidth]{fig_gaze_domain}
\rule{0.7\textwidth}{0.5pt}\\[0.5em]
\textit{[Placeholder: scatter of optical flow vs.\ temporal smoothness with the certified domain $\Omega$ shaded.]}
\caption{Domain-conditioned safe operating region for gaze certification. Inside~$\Omega$ the orientation-and-fixation predicate holds with probability $p_{\text{LB}}\geq 0.82$ at 90\% confidence; outside~$\Omega$ the certificate does not apply. Adapted from \citet{gaze2025reliability}.}
\label{fig:gaze-domain}
\end{figure}



\subsection{Cross-Case Synthesis}
\label{sec:comparison}

The three cases share an architectural shape: a multimodal time-aligned dataset of MR system parameters, head and eye signals, and a self-reported cognitive-state label; a learned predictor; and a verification step layered on top. They differ in which rung of the ladder the verification step sits on, and that difference is what determines the operational claim each case can support. Table~\ref{tab:cross} summarizes the comparison.

\begin{table}[t]
\centering
\caption{Cross-case comparison. Each case study uses a different verification rung (Section~\ref{sec:ladder}), reflecting the operational decision the guarantee is designed to support.}
\label{tab:cross}
\footnotesize
\setlength{\tabcolsep}{3pt}
\begin{tabular}{p{1.7cm}p{2.1cm}p{1.9cm}p{2.3cm}p{2.2cm}}
\hline
\textbf{Case} & \textbf{Predictor} & \textbf{Label} & \textbf{Verification rung} & \textbf{Operational use}\\
\hline
Engagement & Contrastive LSTM + regression & Likert 1--7 self-report & Rung 4: DeepPoly / Prover & Provably safe input intervals\\
Gaze & Time-Aware LSTM & Predicate on gaze output & Rung 2: Clopper--Pearson on predicate & Runtime gate for adaptation\\
Cybersickness (snapshot) & Hierarchical BN & FMS 0--10 (discrete) & Rung 3: Dice on BN posterior & Risk-thresholded intervention\\
Cybersickness (extended) & Random Forest (600 trees) & FMS 0--10 (continuous) & Rung 4: VoRF MILP on ensemble & Pre-session screening\\
\hline
\end{tabular}
\end{table}

The pattern that emerges is not that one rung is best. It is that the right rung is determined by the operational decision the guarantee has to support. A runtime gate (``can the system trust gaze right now?'') wants a runtime-checkable predicate certificate; an intervention trigger (``has cumulative risk crossed a threshold?'') wants a calibrated posterior probability; a screening decision (``is this user safe to admit to a high-intensity session?'') wants a sound bound on the model's worst-case output.

\section{Evaluation, Guarantees, and Limitations}
\label{sec:evaluation}

This section returns to the ladder of Section~\ref{sec:ladder} now that the case studies have grounded each rung in a concrete example, and discusses three things: how performance is measured, what each rung's guarantee licenses an operational system to do, and the threats to validity that bound how far any of these claims can be carried.

\textbf{Performance metrics.} The case studies use complementary metrics, each suited to the prediction problem. The engagement, gaze, and extended-exposure cybersickness models are regression problems and report mean absolute error (MAE), root-mean-squared error (RMSE), and the coefficient of determination ($R^2$). The cybersickness Bayesian network is an ordinal classification problem and reports exact-label accuracy together with within-one-level (``$\pm 1$'') accuracy, the latter being more informative when the cost of a small severity miscalibration is low. Both cybersickness models additionally report sensitivity and specificity at safety thresholds (FMS~$>4$ for moderate symptoms, FMS~$>6$ for severe), which is the right framing for threshold-crossing detection. The gaze case adds a behavioral coupling metric (Pearson $r$ between gaze dwell duration and mission-completion latency) to make the case that the predictor is measuring something operationally meaningful, not just statistically predictable. Calibration matters too. The extended-exposure cybersickness model is well calibrated at low and high perturbation intensity but predicts $P(\text{FMS} > 4) = 0.475$ in the medium-intensity condition where the observed rate is 0.288---a discrepancy the source paper is careful to label as ``an important limitation of the current model rather than evidence of reliable conservatism'' \citep{cybersickness2025extended}.

\textbf{What each rung licenses an operational system to do.} Reading the four rungs from the perspective of an operational system clarifies what they are useful for. Rung~1 supports \emph{average-case decision support}: ``on this kind of cohort, this kind of model performs at this level on average.'' It is appropriate for deciding whether a model is worth deploying at all, but not for individual-decision authority. Rung~2 supports a \emph{runtime gate}: ``inside the certified domain~$\Omega$, the system is licensed to use the gaze-driven adaptation; outside~$\Omega$ it must fall back to mitigation.'' Rung~3 supports a \emph{risk-threshold action}: ``if the posterior $P(\text{FMS} > 4 \mid \text{evidence})$ exceeds~$\varepsilon$, intervene.'' Rung~4 supports \emph{pre-deployment screening}: ``the model is certified never to predict above~4 over this region of admissible inputs, so users whose pre-session features fall in that region can be admitted to the session with the predictor as a safety net.'' These four shapes of guarantee correspond to four different operational architectures and different placements of the verification step in the deployment pipeline. The chapter's central argument is that none of these four substitutes for the others.

\begin{definition}
A \textbf{guarantee} in this chapter is described only at the level supported by the cited work. Specifically: rung~1 is empirical evaluation (average behavior on an exchangeable test set); rung~2 is a statistical certificate on a runtime predicate inside a defined operating region; rung~3 is an exact posterior probability computed under specified evidence on a graphical model; and rung~4 is a sound deterministic bound on the output of a fixed model artifact over a bounded input region.
\end{definition}

\textbf{Threats to validity.} The chapter's case studies are honest about their limits, and the patterns are worth flagging because they are typical of cognitive-state modeling more broadly. \emph{Cohort size and homogeneity}: the four datasets used here have 11, 12, 30, and 32 participants respectively, with limited age and demographic diversity, so generalization claims should be conservative. \emph{Subject-independent versus within-cohort cross-validation}: the extended-exposure cybersickness model explicitly uses band-balanced (not leave-one-subject-out) cross-validation, so the same participant can appear in both training and test folds; the source paper warns that the reported performance should be read as within-cohort rather than truly subject-independent \citep{cybersickness2025extended}. \emph{Self-report dominance}: engagement, FMS, and SSQ are all self-report scales, with the usual recall and demand-characteristic biases; objective physiological correlates exist but are only partially integrated. \emph{Single-testbed validation}: each model is trained and tested on a single MR scenario; transfer to other tasks (medical, educational, gaming) is currently unverified. \emph{Conservatism of deterministic verification}: the engagement DeepPoly result certifies sound but extremely narrow operating intervals; the bound is correct but the certified region is too small to cover most realistic operating points. This is a structural property of high-capacity NN verification, not a failure of the case study. \emph{Calibration in the medium regime}: the extended-exposure cybersickness model is poorly calibrated at medium perturbation intensity; the source paper explicitly does not interpret this as ``conservatism'' but as a model limitation. Each of these threats is a real bound on what the chapter can claim.

\section{Implications and Future Directions}
\label{sec:implications}

This section connects the technical material back to the textbook's broader themes and outlines where the field is likely to move next.

\textbf{Resilience versus mitigation.} The textbook draws a useful distinction between making users \emph{resilient} (so they continue to perform well even when the system degrades) and \emph{mitigating} the degradation when it occurs. The methods in this chapter support both, and they do so in characteristic ways. Certified safe operating regions (rungs~2 and~4) support resilience by construction: the system stays inside a region where the cognitive-state property is known to hold, so the user's attention or comfort is preserved without any reactive intervention. Threshold-crossing predictions (rungs~1 and~3) support mitigation: the system acts when an intervention is needed---rendering stabilization, target enlargement, content simplification, a rest break, or in the limit a session abort. All three case studies provide both shapes. Engagement provides certified safe input intervals plus a runtime engagement estimate that can fire a disengagement alert. Gaze provides a certified domain plus a runtime predicate evaluator that triggers fallback when the certificate no longer applies. The cybersickness BN provides a posterior used for threshold-based intervention, while the cybersickness Random Forest with VoRF provides a screening certificate that selects users into safe regions before they ever start a session.

\textbf{Trust, credibility, and individual differences.} The textbook's recurring question about the role of trust, credibility, and individual differences has a sharp empirical answer in this chapter. In the extended-exposure cybersickness model, baseline susceptibility, hours of sleep, and early-session FMS together account for roughly 62\% of the impurity-based feature importance, while runtime gaze and scene features together account for only about 10\% \citep{cybersickness2025extended}. A similar pattern is visible in the engagement case study, where held-out user variability is non-trivial even for a model that performs well on average. The implication is that an MR cognitive-security system that ignores individual differences in baseline susceptibility---and instead assumes that real-time signals are the dominant source of risk---will be miscalibrated in exactly the cases it most needs to handle. The right architecture combines pre-session screening (where individual differences dominate) with runtime certification (where they are smaller but not zero).

\textbf{Practical application.} For a practitioner, the chapter's framing translates into three concrete deployment patterns. Pre-session, run a screening certificate (rung~4 on a pre-session feature vector) to admit users into safe regions of the operating space. During the session, run a runtime gate (rung~2 on measurable scene statistics) that decides whether a gaze-driven or attention-driven adaptation is licensed. In parallel, run an intervention trigger (rung~3 on a probabilistic posterior) that fires when cumulative risk crosses an application-specific threshold. Each of these uses a different rung because each answers a different operational question.

\textbf{Future directions.} A small number of unresolved problems would meaningfully change the field if solved. First, hybrid symbolic--statistical verification. The deterministic verification result for engagement is sound but conservative; the statistical certification result for gaze is actionable but not deductive; combining the two---formal guarantees on a probabilistic property---would let a system carry a sound claim outside the regions where rung~4 alone is feasible. Second, leave-one-subject-out validation as a methodological default. The within-cohort cross-validation in the extended-exposure cybersickness model is a known limitation, and the field would benefit from systematically reporting the gap between within-cohort and subject-independent performance. Third, multimodal physiological integration. Skin conductance is the only physiological signal carried through the extended-exposure cybersickness model; richer physiology (heart-rate variability, EEG power, respiration) is plausibly available on next-generation MR hardware and would allow better separation of attention, workload, and discomfort. Fourth, larger and more diverse cohorts that span age, gender, and prior MR experience would expand the regions over which the certificates in this chapter actually apply. Fifth, cross-testbed transfer: training and certifying on one MR scenario and evaluating on another is largely untested and is the natural next step.

The next generation of researchers entering this field will inherit a field where the predictive-modeling problems have been largely tamed by deep learning, but where the verification and certification problems are still wide open. The four-rung ladder is a useful map of where the open problems live, and the case studies in this chapter are concrete demonstrations of how each rung looks in practice.

\begin{trailer}{Key Points}
\begin{itemize}
\item Cognitive security in MR is the property that user attention, comfort, and decision-making remain inside operationally safe regions despite measurable degradation of the rendering and sensing pipeline.
\item ``Formal methods'' in this setting is not a single technique. The chapter organizes them as a ladder: empirical evaluation, statistical certification of a runtime predicate, probabilistic verification on a graphical model, and deterministic verification of a fixed model artifact.
\item Each case study exercises a different rung because each answers a different operational question. The gaze case uses Clopper--Pearson certificates; the cybersickness BN uses Dice exact inference; the engagement and cybersickness Random Forest cases use sound verification of a fixed predictor (DeepPoly and VoRF respectively).
\item The right rung for a deployment is determined by the operational decision the guarantee has to support, not by which technique is mathematically strongest.
\item Individual differences and pre-session susceptibility dominate runtime signals in the extended-exposure cybersickness data, which means deployments should combine pre-session screening with runtime certification rather than relying on either alone.
\end{itemize}
\end{trailer}

\begin{trailer}{Questions}
\begin{enumerate}
\item State the four rungs of the empirical-to-formal guarantee ladder introduced in Section~\ref{sec:ladder}, and give one example from this chapter for each rung.
\item Why does the engagement case study report \emph{sound but extremely narrow} certified intervals after DeepPoly verification? What does this tell you about the limits of deterministic neural-network verification on high-capacity models?
\item The gaze case study explicitly does \emph{not} attempt formal NN verification. What does it certify instead, and why is a runtime-checkable certificate the right shape for the gaze-driven-adaptation decision?
\item Compare the two cybersickness models. They predict almost the same thing (FMS severity) but produce different shapes of guarantee. What operational question does each one's guarantee answer?
\item Pearson correlation between gaze dwell duration and rescue latency in the gaze case is $r=0.504$, $p=0.0053$. What does this correlation tell us about the construct validity of using gaze as a proxy for engagement, and what would you need to show to claim a causal relationship?
\item In the extended-exposure cybersickness data, baseline susceptibility and sleep account for roughly 62\% of feature importance while runtime gaze and scene features account for about 10\%. What does this imply for the design of a deployment that wants to support both resilience and mitigation?
\item The cybersickness Random Forest is well calibrated at low and high perturbation intensities but poorly calibrated at medium intensity. Why does this matter operationally, and which rung of the ladder is most affected by miscalibration?
\item A team proposes to deploy an MR cognitive-security system that issues a single ``safe / unsafe'' light to the user based on the output of any of the predictors in this chapter. Critique the proposal in terms of the four rungs.
\end{enumerate}
\end{trailer}

\section{Further Reading}
\begin{enumerate}\leftskip3pt\itemsep6pt
\item \citet{rebenitsch2016cybersickness} for a broad review of cybersickness in applications and visual displays, and \citet{yang2022cybersicknessml} for a more recent systematic review of machine-learning approaches to cybersickness prediction.
\item \citet{plopski2023eye} for a comprehensive survey of gaze interaction and eye tracking in head-worn extended reality.
\item \citet{singh2019deeppoly} for the abstract-interpretation foundation of the DeepPoly verifier used in the engagement case study, and \citet{tornblom2019vorf} for the MILP-based VoRF framework used to certify the cybersickness Random Forest.
\end{enumerate}

\printbibliography
\printindex
\end{document}
